\documentclass{article}
\usepackage[english]{babel}
\usepackage{geometry,amsmath,amssymb,hyperref}
\geometry{letterpaper, portrait}

%%%%%%%%%% Start TeXmacs macros
\newcommand{\cdummy}{\cdot}
\newcommand{\mathd}{\mathrm{d}}
\newcommand{\tmaffiliation}[1]{\\ #1}
\newcommand{\tmem}[1]{{\em #1\/}}
\newcommand{\tmstrong}[1]{\textbf{#1}}
\newcommand{\tmtextit}[1]{\text{{\itshape{#1}}}}
\newenvironment{itemizedot}{\begin{itemize} \renewcommand{\labelitemi}{$\bullet$}\renewcommand{\labelitemii}{$\bullet$}\renewcommand{\labelitemiii}{$\bullet$}\renewcommand{\labelitemiv}{$\bullet$}}{\end{itemize}}
\newtheorem{theorem}{Theorem}
%%%%%%%%%% End TeXmacs macros

\begin{document}

\title{{\tmstrong{An Orthogonal Basis for the Bessel Functions of the First
Kind of Orders 0 and 1}}}

\author{
  Stephen Crowley
  \tmaffiliation{October 10, 2023}
}

\maketitle

\begin{abstract}
  The even-indexed orthonormalized Fourier transforms of the Chebyshev
  polynomials of the first kind form a basis in a reproducing-kernel Hilbert
  space for the Bessel function of the first kind $J_0$ and likewise for the
  odd-indexed functions which form a basis that reproduces $\dot{J}_0 = -
  J_1$. Suprisingly, such a basis for these functions was not known to exist
  before this.
\end{abstract}

{\tableofcontents}

\section{The Type-I Chebyshev Polynomials $T_n (x)$}

Let $T_n$ be the Chebyshev polynomials of the first kind, also said to be of
Type-I, defined by
\begin{equation}
  \begin{array}{ll}
    T_n (x) & =_2 F_1 \left( \begin{array}{cc}
      n, & - n\\
      & \frac{1}{2}
    \end{array} | \frac{1}{2} - \frac{x}{2} \right)\\
    & = \int_{- \infty}^{\infty} e^{ix y}  \hat{T}_n (y) dy\\
    & = \int_{- \infty}^{\infty} e^{ix y}  \frac{i}{y} 
    (\mathrm{e}^{\textrm{-i} y}  F^+_n (y) - \mathrm{e}^{\mathrm{i} y}  (-
    1)^n  F_n^- (y)) dy\\
    & = \int_{- \infty}^{\infty} e^{ix y}  \int_{- \infty}^{\infty} e^{- iy
    z} T_n (z) dz dy
  \end{array}
\end{equation}
where $_2 F_1$ is the (Gauss) hypergeometric function.
{\cite[(13.140)]{ArfkenWeber2005}}

\subsection{The Fourier Transforms $\hat{T}_n (y)$ of $T_n (x)$}

The functions $\hat{T}_n (y)$ are Fourier transforms of $T_n (x)$ defined by
\begin{equation}
  \begin{array}{ll}
    \hat{T}_n (y) & = \int_{- \infty}^{\infty} e^{- ix y} T_n (x) dy = \int_{-
    1}^1 e^{- ix y} T_n (x) dx\\
    & = \int_{- \infty}^{\infty} e^{- ix y} _2 F_1 \left( \begin{array}{cc}
      n, & - n\\
      & \frac{1}{2}
    \end{array} | \frac{1}{2} - \frac{x}{2} \right)   dx\\
    & = \frac{i (e^{- i y} F^+_n (y) - e^{i (\pi n + y)} F_n^- (y))}{y} 
  \end{array}
\end{equation}
where
\begin{equation}
  F_n^{\pm} (y) =_3 F_1 \left( \begin{array}{ccc}
    1, & n, & - n\\
    &  & \frac{1}{2}
  \end{array} | \frac{\pm i y}{2} \right)
\end{equation}
The $L^1$ norm of $\hat{T}_n (y)$ is
\begin{equation}
  | \hat{T}_n | = \sqrt{\frac{4 (- 1)^n \pi - (2 n^2 - 1)}{4 n^2 - 1}}
\end{equation}
then define the normalized Fourier transforms of $T_n (x)$ with
\begin{equation}
  \begin{array}{ll}
    Y_n (y) & = \frac{\hat{T}_n (y)}{| \hat{T}_n |}\\
    & \\
    & = \frac{i}{y} \frac{e^{- i y} F^+_n (y) - e^{i (\pi n + y)} F_n^-
    (y)}{\sqrt{\frac{4 (- 1)^n \pi - (2 n^2 - 1)}{4 n^2 - 1}}}
  \end{array}
\end{equation}
\subsection{The Orthogonalization $Y^{\perp}_n (y)$ of $Y _n (y)$ Via The
Gram-Schmidt Process}

Apply the Gram-Schmidt process to the normalized Fourier transforms of the
Type $I$ Chebyshev polynomials $Y_n (y)$ to get $Y^{\perp}_n (y)$ \
\begin{equation}
  Y^{\perp}_n (y) = Y_n (y) - \sum_{m = 1}^{n - 1} \frac{\langle Y_n (y),
  Y^{\perp}_m (y) \rangle}{\langle Y^{\perp}_m (y), Y^{\perp}_m (y) \rangle}
  Y^{\perp}_m (y)
\end{equation}
then the limits of $Y^{\perp}_n (y)$ at y=0 are equal to
\begin{equation}
  \lim_{y \rightarrow 0} Y^{\perp}_n (y) = \left\{\begin{array}{ll}
    \frac{1}{\sqrt{\pi}} & n = 0\\
    0 & n \neq 0
  \end{array}\right.
\end{equation}
Let
\begin{equation}
  A_{k, n} = - (- 1)^{n + \binom{k}{2}} 2^{2 n - 1 - k} \binom{k + 1}{k - 2 n
  + 1} (k + 2)_{k - 2 n + 1}
\end{equation}
and
\begin{equation}
  B_{k, n} = \cdots
\end{equation}
then we $Y^{\perp}_n (y)$ defined by
\begin{equation}
  \ldots .
\end{equation}
where the first several row-vectors are shown in Table \ref{A}.

\begin{table}[h]
  $\left[ \begin{array}{cccccccc}
    - 3 & 1 & 0 & 0 & 0 & 0 & 0 & \ldots\\
    15 & - 6 & 0 & 0 & 0 & 0 & 0 & \ldots\\
    105 & - 45 & 1 & 0 & 0 & 0 & 0 & \ldots\\
    - 945 & 420 & - 15 & 0 & 0 & 0 & 0 & \ldots\\
    - 10395 & 4725 & - 210 & 1 & 0 & 0 & 0 & \ldots\\
    135135 & - 62370 & 3150 & - 28 & 0 & 0 & 0 & \ldots\\
    2027025 & - 945945 & 51975 & - 630 & 1 & 0 & 0 & \ldots\\
    - 34459425 & 16216200 & - 945945 & 13860 & - 45 & 0 & 0 & \ldots\\
    - 654729075 & 310134825 & - 18918900 & 315315 & - 1485 & 1 & 0 & \ldots\\
    13749310575 & - 6547290750 & 413513100 & - 7567560 & 45045 & - 66 & 0 &
    \ldots\\
    \ldots & \ldots & \ldots & \ldots & \ldots & \ldots & \ldots & \ldots
  \end{array} \right]$
  \caption{The row-vectors \label{A}$A_{k, n}$}
\end{table}

\

\section{Riesz Representation
Theorem}\label{riesz-representation-theorem}\label{riesz-representation-theorem}

The Riesz Representation Theorem serves as a key result in functional
analysis, establishing a connection between linear functionals and measures or
inner products, depending on the context. The theorem has multiple versions,
and two notable ones are as
follows:\label{riesz-representation-theorem-for-hilbert-spaces}\label{riesz-representation-theorem-for-hilbert-spaces}

\begin{theorem}
  {\tmstrong{Riesz Representation Theorem for Hilbert Spaces}}
  
  Let $H$ be a Hilbert space over $\mathbb{C}$ or $\mathbb{R}$, and let $f : H
  \rightarrow \mathbb{C}$ (or $f : H \rightarrow \mathbb{R}$) be a continuous
  linear functional. Then there exists a unique vector $y \in H$ such that for
  every $x \in H$
  \begin{equation}
    f (x) = \langle x, y \rangle
  \end{equation}
  Here, $\langle \cdummy, \cdummy \rangle$ denotes the inner product in $H$.
\end{theorem}

\begin{theorem}
  {\tmstrong{Riesz Representation Theorem for Measures}}
  
  Let $X$ be a locally compact Hausdorff space, and let $C_c (X)$ be the space
  of all continuous functions $f : X \rightarrow \mathbb{C}$ (or $\mathbb{R}$)
  with compact support. If $\Lambda : C_c (X) \rightarrow \mathbb{C}$ (or
  $\mathbb{R}$) is a continuous linear functional, then there exists a unique
  regular complex Borel measure $\mu$ (or real Borel measure, depending on the
  field) on $X$ such that
  \begin{equation}
    \Lambda (f) = \int_X fd \mu
  \end{equation}
  for all $f \in C_c (X)$.
\end{theorem}

The Hilbert spaces version implies that every continuous linear functional can
be represented using the inner product with a specific vector in that space
and the measure version indicates that continuous linear functionals on $C_c
(X)$ can be represented as integrals against a unique measure.

\subsection{Application to Legendre Polynomials}

Considering the space $L^2 ([- 1, 1])$ of Lebesgue square-integrable functions
on the interval $[- 1, 1]$ with the inner product defined as
\begin{equation}
  \langle f, g \rangle = \int_{- 1}^1 f (x) g (x)  \hspace{0.17em} dx
\end{equation}
then it can be shown that the Legendre polynomials $P_n (x)$ are orthogonal to
this inner product:
\begin{equation}
  \langle P_m, P_n \rangle = \int_{- 1}^1 P_m (x) P_n (x)  \hspace{0.17em} dx
  = \delta_{mn}
\end{equation}
where $\delta_{mn}$ is the Kronecker delta. Now, let's say we have a
continuous linear functional $f : L^2 ([- 1, 1]) \to \mathbb{R}$. By the Riesz
Representation Theorem, there exists a unique $g \in L^2 ([- 1, 1])$ such that
\begin{equation}
  f (h) = \langle h, g \rangle
\end{equation}
for all $h \in L^2 ([- 1, 1])$. We can expand $g$ in terms of the Legendre
polynomials:
\begin{equation}
  g (x) = \sum_{n = 0}^{\infty} a_n P_n (x)
\end{equation}
Then,
\begin{equation}
  f (h) = \langle h, g \rangle = \langle h, \sum_{n = 0}^{\infty} a_n P_n
  \rangle = \sum_{n = 0}^{\infty} a_n  \langle h, P_n \rangle
\end{equation}
Here, $a_n$ can be found as:
\begin{equation}
  a_n = \langle g, P_n \rangle
\end{equation}
This can be seen as being equivalent to the application of the Karhunen-Loeve
theorem

\begin{theorem}
  {\tmstrong{The Karhunen-Lo{\`e}ve theorem}}
  
  Let $X_t$ be a zero-mean, square-integrable stochastic process with
  \href{AutocovarianceFunction}{autocovariance function} $R (s, t)$. The
  Karhunen-Lo{\`e}ve theorem states that there exists an orthonormal set of
  functions $\{\phi_k (t)\}$ such that any random variable $X_t$ can be
  represented as an infinite series:
  \begin{equation}
    X_t = \sum_{k = 1}^{\infty} Z_k \phi_k (t)
  \end{equation}
  Here, the $Z_k$ are uncorrelated random variables defined by
  \begin{equation}
    Z_k = \int X_t \phi_k (t) dt
  \end{equation}
  and the functions $\phi_k (t)$ are the eigenfunctions of the autocovariance
  function, which means that
  
  {\tmem{\begin{itemizedot}
    \item They satisfy the integral equation:
    \begin{equation}
      \int R (s, t) \phi_k (t) dt = \lambda_k \phi_k (s)
    \end{equation}
    where $\lambda_k$ are the eigenvalues of the autocovariance function and
    are always real and nonnegative.
    
    \item The eigenfunctions are orthogonal, i.e.,
    \begin{equation}
      \int \phi_k (t) \phi_j (t) dt = \delta_{kj}
    \end{equation}
    where $\delta_{kj}$ is the Kronecker delta function, which equals to 1 if
    $k = j$ and 0 otherwise
    
    \item The uncorrelated random variables $Z_k$ have zero mean and variances
    equal to the corresponding eigenvalues:
    \begin{equation}
      E [Z_k^2] = \lambda_k \text{}
    \end{equation}
    and
    \begin{equation}
      E [Z_k] = 0 \text{}
    \end{equation}
  \end{itemizedot}}}
\end{theorem}

\section{Stationary Gaussian Processes}

\subsection{Bulinskaya's Theorem: $\mathbb{P} [X (s) = x \cap \dot{X} (s) =
0]$=0}

\begin{theorem}
  {\tmstrong{Bulinskaya's theorem on the non-tagency of level crossing
  functionals of continuously differentiable Gaussian processes}}.
  
  Let $X = \{ X_t : 0 \leqslant s \leqslant t \}$ be an almost surely
  continuously differentiable stochastic process with $1$-dimensional density
  $p_s (u)$ bounded in the $u$ variable $\forall s$. \ {\cite[Theorem
  1.1]{GaussianProcessLevelCrossings}}
\end{theorem}

\section{The Gaussian Process $Z (x)$ Having Covariance Kernel $\pi J_0 (2 \pi
t)$}

Consider a zero-mean Gaussian process $Z (t)$ with covariance function given
by
\begin{equation}
  r (| t - s |) =\mathbb{E} [Z (t) Z (s)] = \pi J_0 (2 \pi | t - s |)
\end{equation}
where $E$ denotes expectation. For a stationary zero-mean Gaussian process,
Rice's formula{\cite[Theorem 2.2]{GaussianProcessLevelCrossings}} provides an
expression for the expected number of zeros in a certain interval. For $Z (t)$
with derivative $Z' (t)$ and given that $Z (t)$ and $Z' (t)$ are jointly
Gaussian, the expected number of zeros in the interval $[0, R]$ is given by:
\begin{equation}
  \mathbb{E} [N (R)] = R \int_{- \infty}^{\infty} | \dot{r (t)} | \phi (r (t))
  \mathd t
\end{equation}
where $\phi$ is the standard Gaussian density function and $r (t)$ is the
normalized covariance function. Now compare the result from the Gaussian
process approach with the known result regarding the number of zeros of $J_0$.
{\cite{esgpz}} The power spectral density of the process is given by the
Fourier transform of the covariance kernel $r$
\[ \dot{\rho} (t) = \int_{- \infty}^{\infty} e^{- 2 \pi i s t} r (s, t) \mathd
   s \]
and conversely the covariance kernel is the inverse Fourier transform of the
spectral density
\begin{equation}
  r (h) = \int_{- 1}^1 e^{- \lambda h} \mathd \rho (\lambda)
\end{equation}
where $\rho (s)$ is the spectral distribution function
\begin{equation}
  \int_{- 1}^t \dot{\rho} (s) \mathd s = \int_{- 1}^t \frac{1}{\sqrt{1 - s^2}}
  \mathd s = \int_{- 1}^t \mathd \rho (s) = \frac{\arcsin (s) +
  \frac{\pi}{2}}{\pi}
\end{equation}
and when the integral is written with respect to the function $\rho (s)$
rather than $s$ it symbolizes the Lebesgue-Stieltjes integral with respect to
the spectral distribution function which is absolutely continuous in this case
and therefore the spectral distribution is simply the integral $\rho (t) =
\int \dot{\rho} (t)$ which is also sometimes called the integrated spectrum.
{\cite{yaglom1987correlation}}{\cite{rao2014stochastic}}

\begin{thebibliography}{1}
  \bibitem[1]{ArfkenWeber2005}G.~Arfken  and  H.~Weber.
  {\newblock}\tmtextit{Mathematical Methods for Physicists}.
  {\newblock}Elsevier AP, Boston, 6th  edition, 2005.{\newblock}
  
  \bibitem[2]{GaussianProcessLevelCrossings}Marie~F.~Kratz. {\newblock}Level
  crossings and other level functionals of stationary Gaussian processes.
  {\newblock}\tmtextit{Probability Surveys}, 3:230--288, 2006.{\newblock}
  
  \bibitem[3]{rao2014stochastic}M.M.~Rao. {\newblock}\tmtextit{Stochastic
  Processes - Inference Theory}. {\newblock}Springer Monographs in
  Mathematics. Springer International Publishing, 2014.{\newblock}
  
  \bibitem[4]{yaglom1987correlation}A.M.~Yaglom.
  {\newblock}\tmtextit{Correlation Theory of Stationary and Related Random
  Functions: Volume I: Basic Results}. {\newblock}Applied Probability.
  Springer New York, 1987.{\newblock}
  
  \bibitem[5]{esgpz}N.~Donald Ylvisaker. {\newblock}The Expected Number of
  Zeros of a Stationary Gaussian Process. {\newblock}\tmtextit{The Annals of
  Mathematical Statistics}, 36(3):1043--1046, 1965.{\newblock}
\end{thebibliography}

\

\end{document}
